%%%%%%%%%%%%%%%%%%%%%%%%%%%%%%%%%%%%%%%%%%%%%%%%%%%%%%%%%%%%%%%%%%%%%%
% Lớp 10 | Xác suất | Xác suất của biến cố
% Dạng 0D0?2-4: Liên quan việc chọn người
% Số câu: 35  (sinh tự động bởi tach_cau.py; từ nay thầy sửa trực tiếp file này)
%%%%%%%%%%%%%%%%%%%%%%%%%%%%%%%%%%%%%%%%%%%%%%%%%%%%%%%%%%%%%%%%%%%%%%

%%%%%%%%%%%%%%%% TRẮC NGHIỆM 4 PHƯƠNG ÁN %%%%%%%%%%%%%%%%

% nguồn: MaHoaNH_TracNghiem_Lop10
\begin{ex}%[0D0H2-4]%[ma:632f251a]
	Trong một căn nhà có hai phòng nghỉ ngơi $X$, $Y$ (mỗi phòng có thể chứa được tối đa $3$ người). Ba bạn Sơn, Hải, Văn mỗi bạn  chọn ngẫu nhiên một phòng để nghỉ. Tính xác suất của biến cố \lq\lq cả ba bạn vào cùng phòng\rq\rq.
	 \choice
		 {$\dfrac{1}{5}$}
		{$\dfrac{1}{3}$} 
		{\True  $\dfrac{1}{4}$}
		 {$\mathrm{P}=1$}
	\loigiai{ 
		Các kết quả có thể là $XXX$, $XXY$, $XYX$, $XYY$,  $YXX$, $YXY$, $YYX$, $YYY$.\\
		Do đó số phần tử  của không gian mẫu là $n(\Omega)=8$.\\
		Gọi $H$  là biến cố \lq\lq Cả ba bạn vào cùng phòng\rq\rq. Khi đó $H=\left\{XXX; YYY\right\}$$\Rightarrow n(H)=2\Rightarrow \mathrm{P}(H)=\dfrac{1}{4}$.
	 }
\end{ex}

% nguồn: MaHoaNH_TracNghiem_Lop10
\begin{ex}%[0D0H2-4]%[ma:776c004b]
	Một tổ trong lớp $10A$ có $4$ bạn nam  và $6$ bạn nữ. Giáo viên chọn ngẫu nhiên hai bạn trong tổ để tham gia đội làm báo của lớp. Xác suất để hai bạn được chọn có một bạn nam và một bạn nữ là?
	 \choice
		 {$\dfrac{2}{15}$}
		{$\dfrac{6}{25}$} 
		{$\dfrac{8}{25}$} 
		{\True  $\dfrac{8}{15}$}
	\loigiai{ 
		Số phần tử của không gian mẫu là $n(\Omega) =  \mathrm{C}_{10}^{2} = 45$.\\
		Gọi $A\colon$ \lq\lq Hai bạn được chọn có một bạn nam và  một bạn nữ \rq\rq.\\
		 Khi đó $n(A) = \mathrm{C}_{4}^{1}\cdot \mathrm{C}_{6}^{1} = 4\cdot 6 = 24$.\\
		 Xác suất để hai bạn được chọn có một bạn nam và một bạn nữ là
		$$ \mathrm{P}(A)  = \dfrac{n(A)}{n(\Omega)} = \dfrac{24}{45} = \dfrac{8}{15}. $$
	} 
\end{ex}

% nguồn: MaHoaNH_TracNghiem_Lop10
\begin{ex}%[0D0H2-4]%[ma:872d4ae3]
	Ban chấp hành Đoàn thanh niên của nhà trường cần lập $4$ đội cờ đỏ để chấm thi đua, mỗi đội $3$ người từ $12$ học sinh gồm $5$ học sinh lớp $12$, $4$ học sinh  lớp $11$, $3$ học sinh lớp $10$. Xác suất để đội nào cũng có học sinh lớp $12$ và học sinh lớp $11$ là
	 \choice
		 {\True $\dfrac{36}{385}$}
		 {$\dfrac{6}{385}$}
		{$\dfrac{3}{770}$} 
		 {$\dfrac{9}{385}$}
	\loigiai{ 
		Số cách chia  $12$ học sinh thành $4$ đội là
		 $n\left(\Omega\right)=\mathrm{C}_{12}^3\cdot\mathrm{C}_9^3\cdot\mathrm{C}_6^3\cdot\mathrm{C}_3^3$.\\
		Gọi  $A$ là biến cố \lq\lq Mỗi đội luôn có học sinh lớp $12$ và học sinh lớp $11$\rq\rq.\\
		Xếp vào mỗi đội một học sinh lớp $11$ ta có $4!$  cách.\\
		Xếp $5$ học  sinh lớp $12$ vào $4$ đội thì sẽ có $1$ đội có $2$ học sinh lớp $12$.\\
		Chọn đội có $2$ học sinh lớp $12$ có $4$ cách, chọn $2$ học sinh lớp $12$ có $ \mathrm{C}_5^2$ cách, xếp $3$ học sinh lớp  $12$ còn lại có $3!$ cách.\\
		Xếp $3$  học sinh lớp $10$ có $3!$ cách.\\
		Ta có $ n(A)=4!\cdot4\cdot\mathrm{C}_5^2\cdot3!\cdot3!$.\\ 
		Xác suất của biến cố $A$  là $ \mathrm{P}(A)=\dfrac{4!\cdot4\cdot\mathrm{C}_5^2\cdot3!\cdot3!}{\mathrm{C}_{12}^3\cdot\mathrm{C}_9^3\cdot\mathrm{C}_6^3\cdot\mathrm{C}_3^3}=\dfrac{36}{385}$.
	 }
\end{ex}

% nguồn: MaHoaNH_TracNghiem_Lop10
\begin{ex}%[0D0H2-4]%[ma:243fe18d]
	Đội  tuyển học sinh giỏi Toán $12$ của trường $X$ gồm $8$ học sinh, trong đó có $5$ học sinh nam. Chọn ngẫu nhiên $5$ học sinh đi thi học sinh giỏi cấp tỉnh. Xác suất để $5$ học sinh được chọn đi thi có cả nam và cả nữ học sinh nam nhiều hơn học sinh nữ là
	 \choice
		 {$\dfrac{11}{56}$}
		{\True  $\dfrac{45}{56}$}
		 {$\dfrac{46}{56}$}
		{$\dfrac{55}{56}$} 
	\loigiai{ 
		Số phần tử  của không gian mẫu là $ n\left(\Omega\right)=C_8^5=56$.\\
		Gọi $ A$ là biến  cố \lq\lq $5$ học sinh được chọn đi thi có cả nam và nữ và học sinh nam nhiều hơn học sinh nữ\rq\rq.\\
		Xét  các trường hợp xảy ra của biến cố $A$.
		\begin{itemize}
			\item {\bf Trường hợp 1:} $5$ học sinh được chọn  gồm $4$ nam và $1$ nữ.\\
			Số cách chọn  là $ \mathrm{C}_5^4\cdot\mathrm{C}_3^1=15$.
			\item {\bf Trường hợp 2:} $5$ học sinh được  chọn gồm $3$ nam và $2$ nữ.\\
			Số cách chọn  là $\mathrm{C}_5^3\cdot\mathrm{C}_3^2=30$.
		\end{itemize}
		Do  đó $n(A)=15+30=45$.\\
		Suy ra  $\mathrm{P}(A)=\dfrac{n(A)}{n\left(\Omega\right)}=\dfrac{45}{56}$.
	} 
\end{ex}

% nguồn: MaHoaNH_TracNghiem_Lop10
\begin{ex}%[0D0H2-4]%[ma:ac67d971]
	Một đội gồm $5$ nam và $8$ nữ. Lập một nhóm gồm $4$ người hát tốp ca. Xác  suất để trong $4$ người được chọn có ít nhất $3$ nữ là
	\choice 
		{\True $\dfrac{70}{143}$} 
		 {$\dfrac{73}{143}$}
		{$\dfrac{56}{143}$} 
		{$\dfrac{87}{143}$} 
	\loigiai{ 
		Chọn tùy ý $ 4$ người từ  $ 13$ người.\\
		Suy ra số phần tử của không gian mẫu là $n\left(\Omega\right)=\mathrm{C}_{13}^4=715$ .\\ 
		Gọi $A$ là biến cố \lq\lq $4$ người được chọn  có ít nhất $3$ nữ\rq\rq.\\
		Xét các trường  hợp xảy ra của biến cố $A$.
		\begin{itemize}
			\item {\bf Trường hợp 1:}  Chọn $3$ nữ và $1$ nam.\\
			Số cách chọn  là $ \mathrm{C}_8^3\cdot\mathrm{C}_5^1=280$.
			\item  {\bf Trường hợp 2:} Chọn cả $4$ nữ.\\
			Số cách chọn là  $\mathrm{C}_8^4=70$.
		\end{itemize}
		 Do đó $n(A)=280+70=350$.\\
		Suy  ra $\mathrm{P}(A)=\dfrac{n(A)}{n\left(\Omega\right)}=\dfrac{350}{715}=\dfrac{70}{143}$.
	} 
\end{ex}


%%%%%%%%%%%%%%%% ĐÚNG/SAI %%%%%%%%%%%%%%%%

% nguồn: MaHoaNH_DungSai_Lop10
\begin{ex}%[0D0H2-4]%[ma:4412ba78]
	Xếp $3$ học sinh nam và $3$ học sinh nữ vào một bàn dài  có $6$ ghế. Khi đó
	 \choiceTF
		{Số  phần tử của không gian mẫu là $n(\Omega)=3!3!$}
		{\True Số cách sắp xếp $3$ học sinh nam và $3$ học sinh nữ vào một bàn dài có $6$ ghế sao $3$ học sinh  nam luôn ngồi cạnh nhau là $3!\cdot 4!$}
		{Xác suất để các học sinh nam luôn ngồi cạnh nhau  là $\dfrac{1}{30}$}
		{Xác suất để các học sinh nam và nữ ngồi xen kẽ nhau là $\dfrac{3!\cdot  3!}{6!}$}
	\loigiai{\begin{itemchoice}
		\itemch Số cách xếp $3$ học sinh nam và $3$ học sinh nữ vào một bàn dài có $6$ ghế là hoán vị của $6$ phần tử. Do đó  $n(\Omega)=6!$.
		\itemch Ta coi $3$ học sinh nam luôn ngồi cạnh nhau là $1$ vị trí, số cách sắp xếp $6$ học sinh sao cho  $3$ học sinh nam ngồi cạnh nhau là $3!\cdot 4!$ cách.
		\itemch Xác suất để  các học sinh nam luôn ngồi cạnh nhau là $\dfrac{3!\cdot 4!}{6!}=\dfrac{1}{5}$.
		\itemch Số cách sắp xếp  các em nam và nữ ngồi xen kẽ nhau là $2\cdot 3!\cdot 3!$.\\
		Xác suất để các học sinh nam và nữ ngồi xen kẽ  nhau là $\dfrac{2\cdot 3!\cdot 3!}{6!}$.
	\end{itemchoice}
} 
\end{ex}

% nguồn: MaHoaNH_DungSai_Lop10
\begin{ex}%[0D0H2-4]%[ma:df2bd0e1]
	Ban chấp hành đoàn trường gồm $17$ thành viên trong đó có $8$ đoàn viên nam và $9$ đoàn viên nữ. Chọn  ngẫu nhiên $2$ thành viên tham gia trực tuần.
	\choiceTF 
		{Số cách chọn $2$ thành viên là  $72$}
		{\True Số cách chọn hai  thành viên để một người theo dõi sĩ số, một người phát bản tin hàng ngày là $\mathrm{A}_{17}^2$}
		{Biến cố \lq\lq Hai người được chọn đều là đoàn viên nam\rq\rq có số các kết quả thuận lợi là $56$} 
		{\True Xác suất của biến cố $A\colon $\lq\lq Có ít nhất một thành  viên được chọn là đoàn viên nam\rq\rq bằng $\dfrac{25}{34}$}
	 \loigiai{
		Chọn ngẫu nhiên $2$ thành viên từ Ban chấp  hành đoàn trường gồm $17$ thành viên trong đó có $8$ đoàn viên nam và $9$ đoàn viên nữ.
		\begin{itemchoice}
			\itemch Số  cách chọn là $\mathrm{C}_{17}^2 = 136$ cách, suy ra số phần tử của không gian mẫu là $n(\Omega)=136$.
			\itemch Mỗi cách chọn $2$ thành viên để giao hai nhiệm  vụ một người theo dõi sĩ số, một người phát bản tin hàng ngày là một chỉnh hợp chập $2$ của $17$, do đó có $\mathrm{A}_{17}^2$ cách chọn.
			\itemch Mỗi kết quả thuận lợi của biến cố  \lq\lq Hai người được chọn đều là đoàn viên nam\rq\rq là một tổ hợp chập $2$ của $8$, do đó số các kết quả thuận lợi của biến cố là $\mathrm{C}_8^2 = 28$.
			\itemch Ta  có biến cố $\overline{A}\colon$\lq\lq Cả hai thành viên được chọn đều là đoàn viên nữ\rq\rq.\\
			Khi đó $n(\overline{A}) =  \mathrm{C}_9^2=36$, $\mathrm{P}(\overline{A}) =\dfrac{n(\overline{A})}{n(\Omega)} = \dfrac{36}{136}=\dfrac{9}{34}$.\\
			Vậy xác suất của biến cố $A$ là $\mathrm{P}(A)=1-\mathrm{P} (\overline{A}) = 1-\dfrac{9}{34} = \dfrac{25}{34}$. 
		\end{itemchoice}
	 }
\end{ex}

% nguồn: MaHoaNH_DungSai_Lop10
\begin{ex}%[0D0H2-4]%[ma:4784fb2c]
	 Lớp $10$A có $20$ học sinh nam và $15$ học sinh nữ. Bí thư lớp cần chọn ngẫu nhiên $2$ bạn để đi dự đại hội chi đoàn mẫu. Khi đó
	\choiceTF 
		{\True Không gian mẫu  của phép chọn trên có số phần tử là $595$}
		{\True Xác suất chọn được $2$ bạn nữ  là $\dfrac{3}{17}$}
		{\True Xác suất chọn được $2$ bạn nam là $\dfrac{38}{119}$} 
		{Xác  suất chọn được $1$ bạn nam và một bạn nữ là $\dfrac{1}{17}$}
	 \loigiai{
		\begin{itemchoice}
			\itemch Số phần tử của không gian  mẫu là $\mathrm{C}_{35}^2=595$.
			\itemch Gọi $A$ là biến cố \lq\lq  Chọn được $2$ bạn nữ\rq\rq, khi đó $n(A)=\mathrm{C}_{15}^2$.\\
			Vậy  $\mathrm{P}(A)=\dfrac{\mathrm{C}_{15}^2}{\mathrm{C}_{35}^2}=\dfrac{3}{17}$
			\itemch Gọi $B$ là biến cố \lq\lq Chọn được $2$ bạn nam\rq\rq,  khi đó $n(B)=\mathrm{C}_{20}^2$. \\
			 Vậy $\mathrm{P}(B)=\dfrac{\mathrm{C}_{20}^2}{\mathrm{C}_{35}^2}=\dfrac{38}{119}$.
			\itemch Gọi $C$ là biến cố \lq\lq Chọn được $1$  bạn nam và $1$ bạn nữ\rq\rq, khi đó $n(C)=\mathrm{C}_{20}^1\cdot \mathrm{C}_{15}^1$.\\
			 Vậy $\mathrm{P}(C)=\dfrac{\mathrm{C}_{20}^1\mathrm{C}_{15}^1}{\mathrm{C}_{35}^2}=\dfrac{60}{119}$.
		\end{itemchoice}
	} 
\end{ex}

% nguồn: MaHoaNH_DungSai_Lop10
\begin{ex}%[0D0H2-4]%[ma:b9f31d9f]
	Một nhóm gồm $6$  học sinh nữ và $4$ học sinh nam.
	\choiceTF 
		{Chọn ngẫu nhiên $3$  học sinh từ nhóm đó thì có $10$ cách}
		{Chọn $2$ học  sinh nam và $1$ học sinh nữ từ nhóm đó thì có $12$ cách}
		{\True Chọn ngẫu nhiên $3$ học sinh từ nhóm đó, xác suất để chọn được $2$ học sinh nữ và $1$ học sinh nam bằng $\dfrac{1}{2}$} 
		{\True Chọn ngẫu nhiên $3$ học sinh từ nhóm đó, xác suất để $3$ học sinh được chọn đều là  nam bằng $\dfrac{1}{30}$ }
	\loigiai{ 
		\begin{itemchoice}
			\itemch Sai.  Số cách chọn ngẫu nhiên $3$ học sinh là $\mathrm{C}_{10}^3=120$ cách.
			\itemch Sai. Số cách chọn $2$  nam và $1$ nữ là $\mathrm{C}_4^2 \cdot \mathrm{C}_6^1=36$ cách.
			\itemch Đúng. Số phần tử  không gian mẫu $n\left(\Omega\right)= \mathrm{C}_{10}^3=120$.\\
			Gọi $A$ là biến  cố \lq\lq Chọn được $3$ học sinh gồm $2$ học sinh nữ và $1$ học sinh nam\rq\rq.\\
			Số kết quả thuận lợi của biến  cố $A$ là $n(A)=\mathrm{C}_{6}^2\cdot\mathrm{C}_{4}^1=60$.\\
			Vậy xác suất cần  tìm là $P(A)=\dfrac{n(A)}{n(\Omega)}=\dfrac{1}{2}$.
			\itemch Đúng. Số phần tử không gian mẫu:  $n(\Omega)= \mathrm{C}_{10}^3=120$.\\
			Gọi $A$ là biến cố \lq\lq $3$ học sinh được chọn đều là nam  \rq\rq.\\
			Số kết quả thuận lợi của biến cố $A$ là  $n(A)=\mathrm{C}_{4}^3=4$.\\
			Vậy xác suất  cần tìm là $P(A)=\dfrac{n(A)}{n(\Omega)}=\dfrac{1}{30}$.
		\end{itemchoice}
	 }
\end{ex}

% nguồn: MaHoaNH_DungSai_Lop10
\begin{ex}%[0D0H2-4]%[ma:8ab00711]
	Lớp 10A có $40$ học sinh gồm $25$ nữ và $15$ nam. Cần chọn ra  $4$ học sinh của lớp đi lao động. Các khẳng định sau đúng hay sai?
	 \choiceTF
		{\True Xác suất để chọn được  $4$ học sinh trong đó có $3$ nam và $1$ nữ là $\dfrac{175}{1406}$}
		{\True Xác suất để chọn được $4$ học sinh trong đó có cả nam và nữ là $\dfrac{3150}{9139}$} 
		{Xác suất để  chọn được $4$ học sinh trong đó có $1$ nam và $3$ nữ là $\dfrac{1753}{1406}$}
		{Xác suất để chọn  được $4$ học sinh nam là $\dfrac{1265}{9139}$}
	 \loigiai{
		\begin{itemchoice}
			 \itemch Số cách chọn $4$ trong số $40$ học sinh là $\mathrm{C}_{40}^2 = 91390 \Rightarrow \mathrm{n}(\Omega)= 91390 $.\\
			Gọi $A$ là biến cố  \lq \lq $4$ học sinh được chọn trong đó có $3$ nam và $1$ nữ \rq \rq. \\
			Số kết quả thuận lợi cho biến cố $A$ là $\mathrm{C}_{25}^1 \cdot \mathrm{C}_{15}^3 = 11375 \Rightarrow \mathrm{n}(A) =  11375$.\\
			Vậy $\mathrm{P}(A) = \dfrac{\mathrm{n} (A)}{\mathrm{n}(\Omega)}  = \dfrac{11375}{91390}$.
			\itemch  Số cách chọn $4$ trong số $40$ học sinh là $\mathrm{C}_{40}^2 = 91390 \Rightarrow \mathrm{n}(\Omega)= 91390 $.\\
			Gọi $A$ là biến cố \lq \lq $4$ học sinh được chọn có cả nam và nữ \rq  \rq. \\
			Số kết quả thuận lợi cho biến cố $A$ là $\mathrm{C}_{25}^2 \cdot \mathrm{C}_{15}^2 = 31500  \Rightarrow \mathrm{n}(A) = 31500$.\\
			Vậy $\mathrm{P}(A) = \dfrac{\mathrm{n} (A)}{\mathrm{n}(\Omega)} =  \dfrac{3150}{9139}$.
			\itemch  Số cách chọn $4$ trong số $40$ học sinh là $\mathrm{C}_{40}^2 = 91390 \Rightarrow \mathrm{n}(\Omega)= 91390 $.\\
			Gọi $A$ là biến cố \lq \lq $4$ học sinh được  chọn trong đó có $1$ nam và $3$ nữ \rq \rq. \\
			Số kết quả thuận lợi cho biến cố $A$ là $\mathrm{C}_{25}^3 \cdot \mathrm{C}_{15}^1 = 34500  \Rightarrow \mathrm{n}(A) = 34500$.\\
			Vậy $\mathrm{P}(A) = \dfrac{\mathrm{n}  (A)}{\mathrm{n}(\Omega)} = \dfrac{3450}{9139}$.
			\itemch Số cách  chọn $4$ trong số $40$ học sinh là $\mathrm{C}_{40}^2 = 91390 \Rightarrow \mathrm{n}(\Omega)= 91390 $.\\
			Gọi $A$  là biến cố \lq \lq $4$ học sinh được chọn là nam \rq \rq. \\
			Số kết quả thuận lợi cho biến cố $A$ là $\mathrm{C}_{15}^4= 1365 \Rightarrow \mathrm{n}(A) = 1365$.\\ 
			Vậy $\mathrm{P}(A) = \dfrac{\mathrm{n} (A)}{\mathrm{n}(\Omega)}  = \dfrac{21}{1406}$.
		\end{itemchoice}
	} 
\end{ex}

% nguồn: MaHoaNH_DungSai_Lop10
\begin{ex}%[0D0H2-4]%[ma:b1e928da]
	Một tổ trong lớp $10$A có $11$ học sinh,  trong đó có $7$ học sinh nam và $4$ học sinh nữ. Giáo viên chọn ngẫu nhiên $5$ học sinh trong tổ để kiểm tra vở bài tập Toán. Khi đó
	\choiceTF 
		{\True Số phần tử của không  gian mẫu là $462$}
		{\True Xác suất của biến cố  $B$: \lq\lq $5$ học sinh được chọn đều là nam\rq\rq \, là $\dfrac{1}{22}$}
		{ Xác suất của biến cố $C$: \lq\lq Trong $5$ học sinh được chọn có $3$ nam và  $2$ nữ\rq\rq \, là $\dfrac{41}{462}$}
		{\True Xác suất của biến cố $D$:  \lq\lq Trong $5$ học sinh được chọn có ít nhất $2$ nữ\rq\rq \, là $\dfrac{43}{66}$}
	\loigiai{ 
		\begin{itemchoice}
			\itemch Không gian mẫu  là tập tất cả các tập con gồm $5$ học sinh trong $11$ học sinh.\\
			Vậy số phần  tử của không gian mẫu là $n(\Omega)=\mathrm{C}_{11}^5=462$.
			\itemch Chọn $5$ học sinh nam từ $7$ học sinh nam, có $\mathrm{C}_7^5=21$ (cách chọn) $\Rightarrow  n(B)=21$.\\
			Vậy xác suất của biến  cố $B$ là $\mathrm{P}(B)=\dfrac{n(B)}{n(\Omega)}=\dfrac{21}{462}=\dfrac{1}{22}$.
			\itemch Mỗi phần tử của  $C$ được hình thành từ $2$ công đoạn:
			\begin{itemize}
				\item \textbf{Công đoạn 1.} Chọn $3$ học sinh nam từ $7$ học sinh nam, có  $\mathrm{C}_7^3=35$ (cách chọn).
				\item \textbf{Công đoạn 2.} Chọn $2$  học sinh nữ từ $4$ học sinh nữ, có $\mathrm{C}_{4}^2=6$ (cách chọn).
			\end{itemize}
			Theo quy tắc nhân, tập $C$ có $35\cdot 6=210$ (phần tử). \\Vậy $n(C)=210  \Rightarrow \mathrm{P}(C)=\dfrac{210}{462}=\dfrac{5}{11}$.
			\itemch $5$ học sinh được  chọn có ít nhất $2$ học sinh nữ, có $3$ phương án:
			\begin{itemize}
				\item \textbf{Phương án 1.} Trong $5$ học sinh được chọn có $3$ nam và $2$ nữ:  có $210$ (cách chọn).
				\item \textbf{Phương án 2.} Trong $5$ học sinh được chọn có $2$ nam và $3$ nữ:  có $\mathrm{C}_7^2 \cdot \mathrm{C}_{4}^3=84$ (cách chọn).
				\item \textbf{Phương án 3.} Trong $5$ học sinh được chọn có $1$ nam và $4$ nữ: có $\mathrm{C}_7^1 \cdot  \mathrm{C}_{4}^4=7$ (cách chọn).
			\end{itemize}
			Theo  quy tắc cộng, tập $D$ có $210+84+7=301$ (phần tử).\\Vậy $n(D)=300 \Rightarrow \mathrm{P}(D)=\dfrac{301}{462}=\dfrac{43}{66}$.
	\end{itemchoice}}
\end{ex}

% nguồn: MaHoaNH_DungSai_Lop10
\begin{ex}%[0D0H2-4]%[ma:8609484f]
	Một lớp học có $35$ học sinh, trong đó có $20$ học sinh nam và $15$ học sinh nữ. Cần chọn $5$ học sinh trong  lớp để lập đội cờ đỏ.
	\choiceTF 
		{\True Xác suất để chọn được $5$  học sinh này vào một dãy ghế để họp giao ban mỗi tuần là $\dfrac{15}{40579}$}
		{Xác suất để có $1$ học sinh nam làm đội trưởng, $1$ học sinh nam  làm đội phó và có $2$ học sinh nữ là $0{,}27$}
		 {Xác suất để có $1$ học sinh nữ làm đội trưởng, $1$ học sinh nữ làm đội phó và có $3$ học sinh nam là $0{,}57$}
		{\True Xác suất để có $1$ học sinh nam làm đội trưởng, $1$ học sinh nam làm đội phó và có ít  nhất $1$ học sinh nữ là $0{,}67$}
	\loigiai{ 
		\begin{itemchoice}
			\itemch Số cách chọn $5$ trong số $35$ học sinh là $\mathrm{C}_{35}^5 = 324632 \Rightarrow \mathrm{n}(\Omega)= 324632  $.\\
			Gọi $A$ là biến cố \lq \lq chọn được $5$ học sinh này vào một dãy ghế để họp giao ban mỗi tuần \rq \rq $ \Rightarrow \mathrm{n}(A)  =5!=120$.\\
			Vậy $\mathrm{P}(A) =  \dfrac{\mathrm{n} (A)}{\mathrm{n}(\Omega)} = \dfrac{15}{40579}$.
			\itemch Số cách chọn  $5$ trong số $35$ học sinh là $\mathrm{C}_{35}^5 = 324632 \Rightarrow \mathrm{n}(\Omega)= 324632 $.\\
			Gọi $A$ là biến cố \lq \lq chọn được $1$ học sinh nam làm đội trưởng, $1$  học sinh nam làm đội phó và có $2$ học sinh nữ \rq \rq $ \Rightarrow \mathrm{n}(A) =\mathrm{A}_{20}^2 \cdot \mathrm{C}_{15}^3=172900$.\\
			Vậy $\mathrm{P}(A)  = \dfrac{\mathrm{n} (A)}{\mathrm{n}(\Omega)}\approx 0{,}53$.
			\itemch  Số cách chọn $5$ trong số $35$ học sinh là $\mathrm{C}_{35}^5 = 324632 \Rightarrow \mathrm{n}(\Omega)= 324632 $.\\
			Gọi $A$ là biến cố \lq \lq chọn được $1$ học sinh nữ làm đội trưởng,  $1$ học sinh nữ làm đội phó và có $3$ học sinh nam \rq \rq $ \Rightarrow \mathrm{n}(A) =\mathrm{A}_{15}^2 \cdot \mathrm{C}_{20}^3=23900$.\\
			 Vậy $\mathrm{P}(A) = \dfrac{\mathrm{n} (A)}{\mathrm{n}(\Omega)}\approx 0{,}74$.
			\itemch  Số cách chọn $5$ trong số $35$ học sinh là $\mathrm{C}_{35}^5 = 324632 \Rightarrow \mathrm{n}(\Omega)= 324632 $.\\
			Gọi $A$ là biến cố \lq \lq chọn được $1$ học sinh nam làm đội trưởng, $1$ học sinh nam làm đội phó  và có ít nhất $1$ học sinh nữ \rq \rq.
			\begin{itemize}
				\item \textbf{TH1.} có đúng $1$  học sinh nữ: có $\mathrm{A}_{20}^2 \cdot 15=5700$ cách chọn.
				\item \textbf{TH2.} có đúng $2$  học sinh nữ: có $\mathrm{A}_{20}^2 \cdot \mathrm{C}_{15}^2=39900$ cách chọn.
				\item \textbf{TH3.} có đúng $3$ học sinh nữ: có $\mathrm{A}_{20}^2 \cdot \mathrm{C}_{15}^3=172900$  cách chọn.
			\end{itemize}
			$ \Rightarrow \mathrm{n}(A) =5700+39900+172900=218500$.\\ 
			Vậy $\mathrm{P}(A)  = \dfrac{\mathrm{n} (A)}{\mathrm{n}(\Omega)}\approx 0{,}67$.
		\end{itemchoice}
	 }
\end{ex}

% nguồn: MaHoaNH_DungSai_Lop10
\begin{ex}%[0D0H2-4]%[ma:7cf9ed47]
	Một tổ trong lớp $10B$ có $10$ học sinh, trong  đó có $7$ học sinh nam và $3$ học sinh nữ. Giáo viên chọn ngẫu nhiên $5$ học sinh trong tổ để tập văn nghệ cho đợt 26/3. Khi đó
	 \choiceTF
		{Số phần tử của không gian mẫu  là $560$}
		{\True Xác suất của biến cố  $B\colon$\lq\lq $5$ học sinh được chọn đều là nam\rq\rq\, là $\dfrac{1}{12}$}
		{Xác suất của biến cố $C\colon$ \lq\lq Trong $5$ học sinh được chọn có  $3$ nam và $2$ nữ\rq\rq\, là $\dfrac{41}{462}$}
		{\True Xác suất của biến cố $D\colon$ \lq\lq Trong $5$ học sinh được chọn  có ít nhất $2$ nữ\rq\rq\, là $\dfrac{1}{2}$}
	 \loigiai{
		\begin{itemchoice}
			 \itemch
			Chọn ngẫu nhiên $5$ học sinh trong tổ có $10$  hoc sinh.\\
			Số tử  của không gian mẫu là $n\left(\Omega\right)=\mathrm{C}_{10}^5=252\ne 560$.
			\itemch 
			Biến cố $B\colon$\lq\lq $5$ học sinh được chọn đều là  nam\rq\rq.\\
			Số phần tử biến cố $B$ là  $n(B)=\mathrm{C}_{7}^5=21$.\\
			Suy ra xác suất của biến cố $B$  là
			 $\mathrm{P}(B)=\dfrac{21}{252}=\dfrac{1}{12}$.
			\itemch 
			Biến cố $C\colon$ \lq\lq Trong $5$ học sinh được chọn có $3$ nam  và $2$ nữ\rq\rq.\\
			Số phần  tử biến cố $C$ là $n(C)=\mathrm{C}_{7}^3\cdot\mathrm{C}_{3}^2=105$.\\
			 Suy ra xác suất của biến cố $C$ là
			 $\mathrm{P}(C)=\dfrac{105}{252}=\dfrac{5}{12}\ne \dfrac{41}{462}$.
			\itemch 
			Biến cố $D\colon$ \lq\lq Trong $5$ học sinh được chọn có ít  nhất $2$ nữ\rq\rq.\\
			Các trường hợp chọn biến cố $D$ như sau 
			\begin{itemize}
				\item {\bf Trường hợp 1:} Chọn  $2$ học sinh nữ và $3$ học sinh nam.\\
				Số cách  chọn là $\mathrm{C}_3^2\cdot\mathrm{C}_{7}^3=105$.
				\item {\bf Trường hợp 2:} Chọn $3$ học sinh nữ và $2$ học sinh  nam.\\
				Số  cách chọn là $\mathrm{C}_3^3\cdot\mathrm{C}_{7}^2=21$.
			\end{itemize}
			Do đó $n(D)=105+21=126$.\\ 
			Suy ra xác suất  của biến cố $C$ là
			$\mathrm{P}(C)=\dfrac{126}{252}=\dfrac{1}{2}$. 
		\end{itemchoice}
	} 
\end{ex}

% nguồn: MaHoaNH_DungSai_Lop10
\begin{ex}%[0D0H2-4]%[ma:295fc178]
	Lớp $11A$ có $7$ học sinh nữ và $13$ học sinh nam. Cô chủ  nhiệm chọn ra $5$ bạn để tham gia văn nghệ. Khi đó
	 \choiceTF
		{\True Xác suất để chọn được $5$ học sinh nữ là  $\dfrac{21}{15504}$}
		{\True Xác suất để chọn được đúng $3$ học sinh nam là  $\dfrac{\mathrm{C}_{13}^3\cdot \mathrm{C}_7^2}{\mathrm{C}_{20}^5}$}
		{Xác suất để chọn được ít nhất $1$ học sinh nữ  là $\dfrac{429}{5168}$}
		{\True Xác suất để chọn được học sinh nữ nhiều hơn học sinh số nam là  $\dfrac{1603}{7752}$}
	 \loigiai{
		Chọn ra $5$ bạn trong $7$ học sinh nữ và  $13$ học sinh nam tham gia văn nghệ. Số phần tử không gian mẫu là $n\left(\Omega\right)=\mathrm{C}_{20}^5$.
		\begin{itemchoice}
			\itemch 
			Gọi  $A$ là biến cố \lq\lq chọn được $5$ học sinh nữ là\rq\rq.\\
			 Suy ra $n(A)=\mathrm{C}_{7}^5$.\\
			 Vậy xác suất để chọn được $5$ học sinh nữ là là
			 $\mathrm{P}(A)=\dfrac{n(A)}{n\left(\Omega\right)}
			=\dfrac{\mathrm{C}_{7}^5}{\mathrm{C}_{20}^5}=\dfrac{21}{15\,504}$. 
			\itemch 
			Gọi $B$  là biến cố \lq\lq chọn được đúng $3$ học sinh nam\rq\rq.\\
			Suy ra  $n(B)=\mathrm{C}_{13}^3\cdot\mathrm{C}_7^2$.\\
			Vậy xác suất để chọn đúng $3$ học sinh nam  là
			$\mathrm{P}(B)=\dfrac{\mathrm{C}_{13}^3\cdot \mathrm{C}_7^2}{\mathrm{C}_{20}^5}$. 
			 \itemch
			Gọi $C$ là biến cố \lq\lq chọn được  ít nhất $1$ học sinh nữ\rq\rq.\\
			Suy ra  $\overline{C}\colon$ \lq\lq chọn được $5$ học sinh nam\rq\rq.\\
			Suy  ra $n\left(\overline{C}\right)=\mathrm{C}_{13}^5$.
			 Vậy $ \mathrm{P}(C)=1-\mathrm{P}\left(\overline{C}\right)=1-\dfrac{\mathrm{C}_{13}^5}{\mathrm{C}_{20}^5}=\dfrac{4\,739}{5\,168}\ne\dfrac{429}{5168} $.
			 \itemch
			Gọi $D$ là biến cố \lq\lq chọn được học sinh nữ nhiều  hơn số học sinh nam\rq\rq.\\
			Các  trường hợp chọn biến cố $D$ như sau
			\begin{itemize}
				\item {\bf Trường hợp 1:} Chọn $3$ học  sinh nữ và $2$ học sinh nam. Số cách chọn là $\mathrm{C}_7^3\cdot\mathrm{C}_{13}^2$.
				\item {\bf Trường hợp 2:} Chọn  $4$ học sinh nữ và $1$ học sinh nam. Số cách chọn là $\mathrm{C}_7^4\cdot\mathrm{C}_{13}^4$.
				\item  {\bf Trường hợp 3:} Chọn $5$ học sinh nữ. Số cách chọn là $\mathrm{C}_7^5$.
			\end{itemize}
			Suy ra  $n(D)=\mathrm{C}_7^3\cdot\mathrm{C}_{13}^2+\mathrm{C}_7^4\cdot\mathrm{C}_{13}^1+\mathrm{C}_7^5$.\\
			Xác suất để chọn được học sinh nữ nhiều hơn số  học sinh nam là\\
			 $\mathrm{P}(D)=\dfrac{\mathrm{C}_7^3\cdot\mathrm{C}_{13}^2+\mathrm{C}_7^4\cdot\mathrm{C}_{13}^1+\mathrm{C}_7^5}{\mathrm{C}_{20}^5}=\dfrac{1603}{7752}$.
		\end{itemchoice}
	} 
\end{ex}

% nguồn: MaHoaNH_DungSai_Lop10
\begin{ex}%[0D0V2-4]%[ma:005fea74]
	Trong một ban tổ chức gồm $5$ nhân viên đến từ Việt Nam, $7$ nhân viên đến từ Hoa Kỳ và $6$ nhân viên đến từ  Anh.
	 \choiceTF
		{\True Có $210$ cách chọn ra  $3$ nhân viên, mỗi người từ một quốc gia khác nhau}
		{\True Có $\mathrm{C}^2_7$ cách  chọn ra $2$ nhân viên từ Hoa Kỳ}
		{Chọn ngẫu nhiên $2$ nhân viên từ ban, xác suất để chọn được $2$ nhân viên từ hai quốc gia  khác nhau là $\dfrac{203}{272}$}
		{Chọn  ngẫu nhiên $3$ nhân viên từ ban, xác suất để chọn được $3$ nhân viên từ một quốc gia là $\dfrac{35}{816}$}
	\loigiai{ 
		\begin{itemchoice}
			\itemch Số cách chọn ra $3$ nhân viên, mỗi người từ một quốc gia khác nhau là $5\cdot 7\cdot 6=210$. 
			 \itemch Có $\mathrm{C}^2_7$ cách chọn ra $2$ nhân viên từ Hoa Kỳ.
			\itemch Số cách chọn $2$ nhân  viên từ ban: $\mathrm{C}^2_{18}=153$.\\
			Số cách chọn $2$ nhân viên từ hai quốc gia khác nhau  là
			 $$\mathrm{C}^1_5\mathrm{C}^1_7+\mathrm{C}^1_5\mathrm{C}^1_6+\mathrm{C}^1_6\mathrm{C}^1_7=107.$$
			Xác suất để chọn được $2$ nhân viên từ hai quốc gia khác nhau là  $\dfrac{107}{153}$.
			\itemch Số cách chọn $3$ nhân viên từ ban:  $\mathrm{C}^3_{18}=816$.\\
			Số cách chọn $3$  nhân viên từ một quốc gia là $\mathrm{C}^3_5+\mathrm{C}^3_6+\mathrm{C}^3_7=65$.
			Xác suất để chọn được $3$ nhân viên  một quốc gia là $\dfrac{65}{816}$.
	\end{itemchoice}}
\end{ex}

% nguồn: MaHoaNH_DungSai_Lop10
\begin{ex}%[0D0V2-4]%[ma:4e432f58]
	Trong lớp $10A$ có $25$ bạn nam và $21$ bạn nữ. Giáo  viên chọn ngẫu nhiên $3$ bạn trong lớp để làm cán bộ lớp. Xét tính Đúng - Sai của các khẳng định sau.
	 \choiceTF
		{\True Số cách chọn ra $3$ bạn  trong lớp $10A$ là $15180$ (cách)}
		{\True Xác suất của các  biến cố \lq\lq Ba bạn được chọn đều là nam\rq\rq \, bằng $\dfrac{5}{33}$}
		{Xác suất của các biến cố \lq\lq Ba bạn được  chọn đều là nữ\rq\rq \, bằng $\dfrac{133}{1158}$}
		{\True Xác suất của các biến cố \lq\lq Trong ba học sinh được  chọn có hai bạn nam và một bạn nữ\rq\rq \,bằng $\dfrac{105}{253}$}
	\loigiai{ 
		\begin{itemchoice}
			\itemch Số cách chọn ra $3$ bạn trong lớp $10A$ gồm $46$ bạn ($25$ bạn nam và  $21$ bạn nữ) là $\mathrm{C}_{46}^3=15180$ (cách). \\
			\itemch Gọi $A$ là biến cố \lq\lq Ba bạn được chọn đều  là nam \rq\rq.\\
			Ta có, $n(\Omega)=15180$  và $n(A)=\mathrm{C}_{25}^3 =2300$.\\
			Xác suất của biến cố $A$  là $\mathrm{P}(A)=\dfrac{n(A)}{n(\Omega)}=\dfrac{2300}{15180}=\dfrac{5}{33}$.
			\itemch Gọi $B$ là \lq\lq  Ba bạn được chọn đều là nữ\rq\rq. \\
			Số  cách chọn được $3$ bạn nữ từ $21$ bạn nữ là $\mathrm{C}_{21}^3=1330$ (cách).\\
			Suy ra $n(B)=1330$. Xác suất của biến cố  $B$ là $\mathrm{P}(B)=\dfrac{n(B)}{n(\Omega)}=\dfrac{1330}{15180}=\dfrac{133}{1518}$.
			\itemch Gọi $C$ là biến cố \lq\lq Trong ba học sinh được chọn có hai bạn  nam và một bạn nữ\rq\rq.\\
			Số cách chọn được $2$  bạn nam và $1$ bạn nữ là $\mathrm{C}_{25}^2 \cdot \mathrm{C}_{21}^1=6300$ (cách).\\
			Suy  ra $n(C)=6300$. Xác suất của biến cố $C$ là $\mathrm{P}(C)=\dfrac{n(C)}{n(\Omega)}=\dfrac{6300}{15180}=\dfrac{105}{253}$.
		\end{itemchoice}
	} 
\end{ex}

% nguồn: MaHoaNH_DungSai_Lop10
\begin{ex}%[0D0V2-4]%[ma:4e328aff]
	Một  lớp học có $40$ học sinh gồm $25$ nam, trong đó có bạn An và $15$ nữ, trong đó có bạn Huyền. Mỗi khẳng định sau đúng hay sai?
	\choiceTF 
		{\True Có $\mathrm{C}_{40}^5$ cách chọn $5$ học sinh bất kì của  lớp đi lao động}
		{Có $\mathrm{C}_{40}^4$ cách chọn $5$ bạn học sinh của  lớp sao cho trong đó có mặt bạn An }
		{\True Có $\mathrm{C}_{40}^5-\mathrm{C}_{38}^3$ cách chọn $5$ bạn học sinh của lớp trong đó hai bạn An và Huyền không đồng thời  có mặt}
		{\True Chọn ngẫu nhiên $5$ học sinh của lớp đi dự Đại hội. Xác suất chọn được cả nam và nữ đồng thời hai bạn An và Huyền đều có mặt là  $\dfrac{1}{78}$ }
	\loigiai{ 
		\begin{enumerate}
			\item Chọn $5$ học sinh bất kì có $\mathrm{C}_{40}^5$ cách  chọn.\\
			\item Chọn $5$ học sinh sao cho luôn có  mặt An có $\mathrm{C}_{39}^4$.\\
			\item Số cách chọn $5$ bạn học sinh trong đó hai bạn An và Huyền không đồng thời có  mặt là $\mathrm{C}_{40}^5-\mathrm{C}_{38}^3$ cách.
			\item Số cách chọn  $5$ học sinh trong đó có cả nam và nữ đồng thời An và Huyền luôn có mặt gồm các trường hợp sau
			\begin{itemize}
				\item $5$ bạn gồm An, Huyền và $3$ bạn nam có $\mathrm{C}_{24}^3$ cách. 
				\item $5$ bạn gồm An, Huyền và $2$ bạn nam, $1$ bạn  nữ có $\mathrm{C}_{24}^2\cdot \mathrm{C}_{14}^1$ cách.
				\item $5$ bạn gồm An, Huyền và  $1$ bạn nam, $2$ bạn nữ có $\mathrm{C}_{24}^1\cdot \mathrm{C}_{14}^2$ cách.
				\item  $5$ bạn gồm An, Huyền và $3$ bạn nữ có $ \mathrm{C}_{14}^3$ cách.
			\end{itemize}
			Vậy xác suất để chọn được $5$ bạn học  sinh có cả nam và nữ đồng thời hai bạn An và Huyền luôn có mặt là
			\[P=\dfrac{\mathrm{C}_{24}^3+  \mathrm{C}_{24}^2\cdot \mathrm{C}_{14}^1+\mathrm{C}_{24}^1\cdot \mathrm{C}_{14}^2+\mathrm{C}_{14}^3}{\mathrm{C}_{40}^5}=\dfrac{1}{78}\]
		\end{enumerate}
	} 
\end{ex}

% nguồn: MaHoaNH_DungSai_Lop10
\begin{ex}%[0D0V2-4]%[ma:40a70760]
	Một tổ  có $8$ nam và $7$ nữ. Chọn ngẫu nhiên 5 người. Khi đó
	 \choiceTF
		{\True Số phần tử của  không gian mẫu là $3\,003$}
		{\True  Xác suất để có 3 nam và 2 nữ bằng $\dfrac{56}{143}$}
		{Xác suất để có cả nam lẫn nữ mà nam nhiều hơn nữ bằng $\dfrac{23}{429}$} 
		{\True Xác suất để có  ít nhất một nữ bằng $\dfrac{421}{429}$}
	\loigiai{ 
		\begin{itemchoice}
			\itemch Không  gian mẫu: $n\left(\Omega\right)=\mathrm{C}_{15}^5=3\,003$.
			\itemch Gọi A là biến cố: \lq\lq $5$ người được chọn có  $3$ nam và $2$ nữ\rq\rq.\\
			Số cách chọn $5$ người trong đó có $3$ nam, $2$ nữ  là $\mathrm{C}_8^3\cdot \mathrm{C}_7^2.$\\
			$\Rightarrow n(A)=\mathrm{C}_8^3\cdot  \mathrm{C}_7^2.\\
			\Rightarrow  \mathrm{P}(A)=\dfrac{n(A)}{n\left(\Omega\right)}=\dfrac{1\,176}{3\,003}=\dfrac{56}{143}$.
			\itemch Gọi $B$ là biến cố: \lq\lq $5$ người được chọn  có cả nam lẫn nữ mà nam nhiều hơn nữ\rq\rq.\\
			 Trường hợp 1: Chọn $4$ nam và $1$ nữ.\\
			Số cách chọn $5$ người  trong đó có $4$ nam, $1$ nữ là $\mathrm{C}_8^4\cdot \mathrm{C}_7^1$.\\
			 Trường hợp 2: Chọn $3$ nam và $2$ nữ.\\
			Số cách chọn $5$ người trong đó có $3$ nam, $2$ nữ là $\mathrm{C}_8^3\cdot \mathrm{C}_7^2$.\\ 
			$\Rightarrow  n(B)=\mathrm{C}_8^4\cdot \mathrm{C}_7^1+\mathrm{C}_8^3\cdot \mathrm{C}_7^2=1\,666.\\
			\Rightarrow  \mathrm{P}(B)=\dfrac{n(B)}{n\left(\Omega\right)}=\dfrac{1\,666}{3\,003}=\dfrac{238}{429}$.
			\itemch Gọi $C$ là biến  cố: \lq\lq $5$ người được chọn có ít nhất một nữ\rq\rq.\\
			Vậy ta có $\overline{C} \colon$  \lq\lq$5$ người được chọn không có nữ\rq\rq.\\
			$\Rightarrow n\left(\overline{C}\right)=\mathrm{C}_8^5$.\\ 
			Vậy xác suất cần tìm  là $\mathrm{P}(C)=1-\mathrm{P}\left(\overline{C}\right)=1-\dfrac{n\left(\overline{C}\right)}{n\left(\Omega\right)}=1-\dfrac{\mathrm{C}_8^5}{3003}=\dfrac{421}{429}$.
	\end{itemchoice}}
\end{ex}

% nguồn: MaHoaNH_DungSai_Lop10
\begin{ex}%[0D0V2-4]%[ma:f3d3f790]
	Lớp $ 10C5 $ có $ 20 $ học sinh trong đó có $ 12 $ học sinh nữ và $  8 $ học sinh nam. Giáo viên gọi ngẫu nhiên hai học sinh lên bảng làm bài. Chọn khẳng định đúng, khẳng định sai trong các khẳng định sau
	 \choiceTF
		{\True Xác suất để hai học sinh được gọi có đúng một  học sinh nam là $\dfrac{48}{95}$}
		{Xác suất để hai học sinh được  gọi có hai học sinh nữ là $\dfrac{6}{95}$}
		{\True Xác suất để hai học sinh được gọi có ít nhất một học sinh nam  là $\dfrac{62}{95}$}
		{\True Xác suất để hai học sinh nam được gọi  là $\dfrac{14}{95}$}
	\loigiai{ 
		\begin{itemchoice}
			\itemch Đây là khẳng định  đúng vì\\
			Số phần  tử của không gian mẫu là $ n(\Omega)=\mathrm{C}_{20}^2=190$.\\
			Số cách chọn $1$ học sinh  trong đó đúng $1$ học sinh nam là $\mathrm{C}_{12}^1\cdot\mathrm{C}_{8}^1$.\\
			Xác suất để chọn được đúng một học sinh nam là  $\dfrac{\mathrm{C}_{12}^1\cdot\mathrm{C}_{8}^1}{\mathrm{C}_{20}^2}=\dfrac{96}{190}=\dfrac{48}{95}$.
			\itemch Đây là  khẳng định sai vì
			Số cách chọn $2$ học sinh  nữ là $\mathrm{C}_{12}^2$.\\
			 Xác suất để chọn được $2$ học sinh nữ là $\dfrac{\mathrm{C}_{12}^2}{\mathrm{C}_{20}^2}=\dfrac{66}{190}=\dfrac{33}{95}$.
			\itemch Đây là khẳng định  đúng vì\\
			Gọi  $ A $ là biến cố : \lq\lq Trong hai học sinh được gọi có ít nhất $ 1 $ nam\rq\rq.\\
			{\bf Trường hợp 1.} Chọn $ 1 $ nam, $ 1  $ nữ có $ \mathrm{C}_{12}^1\cdot \mathrm{C}_{8}^1 =96$ cách.\\
			{\bf Trường hợp 2.}  Chọn $ 2 $ nam có $ \mathrm{C}_{8}^2=28$ cách.\\
			 Số phần tử của biến cố $ A $ là $ n(A)=96+28=124$.\\
			Suy ra $ \mathrm{P}(A)=\dfrac{n(A)}{n(\Omega)}=\dfrac{124}{190}=  \dfrac{62}{95}$.\\
			\itemch Đây  là khẳng định đúng vì\\
			Gọi $C$  là biến cố : \lq\lq Hai học sinh được gọi là mam \rq\rq.\\
			 Ta có $\mathrm{P}(C)=\dfrac{n(C)}{n(\Omega)}=\dfrac{
			\mathrm{C}_{8}^2}{190}= \dfrac{14}{95}$. 
		\end{itemchoice}
	 }
\end{ex}

% nguồn: MaHoaNH_DungSai_Lop10
\begin{ex}%[0D0V2-4]%[ma:f2f634d5]
	Trong lớp $10C$ có $38$ học sinh gồm $18$ nam và $20$ nữ. Giáo viên chủ nhiệm cần chọn ra $5$ học sinh để thành lập đội văn  nghệ của lớp. Khi đó
	 \choiceTF
		{\True Số phần tử của không gian mẫu là  $n(\Omega)=501\,942$}
		{\True Xác suất của biến cố $B$: \lq\lq Chọn được đội văn nghệ có đúng $3$ học sinh  nữ\rq\rq\ là $n(B)=\dfrac{90}{256}$}
		{Xác suất của biến cố $C$: \lq\lq Chọn được đội văn nghệ có đủ cả  nam và nữ, đồng thời số học sinh nam nhiều hơn số học sinh nữ\rq\rq\ là $n(C)=\dfrac{2\,119}{4\,921}$}
		{\True Xác suất của biến cố $D$: \lq\lq Chọn được đội văn nghệ sao cho có đủ cả nam và nữ, đồng thời số học sinh nam là một  số chẵn\rq\rq\ là $\dfrac{330}{703}$}
	\loigiai{ 
		\begin{itemchoice}
			\itemch Số phần  tử của không gian mẫu là $n(\Omega)=\mathrm{C}_{38}^5=501\,942$.
			\itemch Để chọn được đội văn nghệ gồm $3$ nữ và $2$ nam ta thực hiện qua hai  bước.
			\begin{itemize}
				\item \textbf{Bước 1:} Chọn $3$  nữ từ $20$ nữ có $\mathrm{C}_{20}^3=1\,140$ cách.
				\item  \textbf{Bước 2:} Chọn $2$ nam từ $18$ nam có $\mathrm{C}_{18}^2=153$ cách.
			\end{itemize}
			Áp dụng  quy tắc nhân ta có số phần tử của biến cố $B$ là $n(B)=1\,140\cdot 153=174\,420$.\\
			Vậy xác suất của biến  cố $B$ là $\mathrm{P}(B)=\dfrac{n(B)}{n(\Omega)}=\dfrac{174\,420}{501\,942}=\dfrac{90}{259}$.
			\itemch Để  chọn được đội văn nghệ như yêu cầu đề bài, ta xét hai trường hợp sau
			\begin{itemize}
				\item \textbf{Trường hợp 1.}  Chọn $4$ nam và $1$ nữ, số cách chọn là $\mathrm{C}_{18}^4\cdot \mathrm{C}_{20}^1=61\,200$.
				\item \textbf{Trường hợp 2.}  Chọn $3$ nam và $2$ nữ, số cách chọn là $\mathrm{C}_{18}^3\cdot \mathrm{C}_{20}^2=155\,040$.
			\end{itemize}
			 Áp dụng quy tắc cộng ta có số phần tử của biến cố $C$ là $n(C)=155\,040+61\,200=216\,240$.\\
			Vậy xác suất của biến cố  $C$ là $\mathrm{P}(C)=\dfrac{n(C)}{n(\Omega)}=\dfrac{216\,240}{501\,942}=\dfrac{2\,120}{4\,921}$.
			\itemch Để chọn được đội văn nghệ  như yêu cầu đề bài, ta xét hai trường hợp sau
			\begin{itemize}
				\item \textbf{Trường hợp 1.} Chọn $2$ nam và $3$ nữ, số cách chọn là $\mathrm{C}_{18}^2\cdot \mathrm{C}_{20}^3=174\,420$. 
				\item \textbf{Trường hợp 2.} Chọn $4$ nam  và $1$ nữ, số cách chọn là $\mathrm{C}_{18}^4\cdot \mathrm{C}_{20}^1=61\,200$.
			\end{itemize}
			Áp dụng quy tắc cộng ta có số phần tử  của biến cố $D$ là $n(D)=174\,420+61\,200=235\,620$.\\
			Vậy xác suất  của biến cố $D$ là $\mathrm{P}(D)=\dfrac{235\,620}{501\,942}=\dfrac{330}{703}$.
		\end{itemchoice}
	} 
\end{ex}

% nguồn: MaHoaNH_DungSai_Lop10
\begin{ex}%[0D0V2-4]%[ma:3ca9702d]
	Từ tập hợp gồm $12$ học sinh trong đó  có $6$ học sinh khối $12$, $4$ học sinh khối $11$ và $2$ học sinh khối $10$, người ta chọn hoặc cho học sinh xếp hàng.
	\choiceTF 
		{Có $27$ cách chọn  ra $3$ học sinh cùng một khối}
		{\True Có $48$ cách chọn ra $3$ học sinh ở $3$  khối}
		{\True Có $207\ 360$ cách xếp $12$  học sinh trên thành một hàng ngang sao cho các học sinh cùng khối thì đứng cạnh nhau}
		{ Có $1\ 036\ 800$ cách xếp $12$ học  sinh trên thành một hàng ngang sao cho không có $2$ học sinh cùng khối nào đứng cạnh nhau}
	 \loigiai{
		\begin{itemchoice}
			\itemch {\bf Sai.} Để chọn học sinh cùng một khối  tham gia đội xung kích, ta có các phương án
			\begin{enumerate}
				\item Phương án 1: $3$ học sinh cùng khối  12, có $\mathrm{C}_6^3$ cách.
				\item Phương án 2: $3$ học  sinh cùng khối 11, có $\mathrm{C}_4^3$ cách.
			\end{enumerate}
			Suy ra số cách chọn $3$ học sinh cùng khối  là $\mathrm{C}_6^3+\mathrm{C}_4^3=24$.
			\itemch {\bf Đúng.} Chọn  $1$ học sinh khối $12$ có $6$ cách.\\
			Chọn $1$  học sinh khối $11$ có $4$ cách.\\
			Chọn $1$ học sinh khối $10$ có $2$  cách.\\
			Theo quy tắc nhân, số cách là $6\cdot 4\cdot 2=48$  (cách).
			\itemch  {\bf Đúng.} Sắp xếp $6$ học sinh khối 12 với nhau, có $6!$ cách.\\
			Sắp xếp $4$ học sinh  khối 11 với nhau, có $4!$ cách.\\
			Sắp xếp  $2$ học sinh khối 10 với nhau, có $2!$ cách.\\
			Sắp xếp $3$  khối $10, 11, 12$ với nhau, có $3!$ cách.\\
			Theo quy tắc nhân, số cách sắp  xếp là $6!\cdot 4!\cdot 2!\cdot 3!=207\ 360$ (cách).
			\itemch  {\bf Sai.} Trước hết, xếp $6$ học sinh khối 12 thành hàng ngang, có $6!$ cách.
			\begin{center}
				\scalebox{1.5}{\color{teal}\circled{$1$}}\ \scalebox{2}{\textcolor{red}{\fbox{\bf $12$}}}\  \scalebox{1.5}{\color{teal}\circled{$2$}}\ \scalebox{2}{\textcolor{red}{\fbox{\bf $12$}}}\ \scalebox{1.5}{\color{teal}\circled{$3$}}\ \scalebox{2}{\textcolor{red}{\fbox{\bf $12$}}}\ \scalebox{1.5}{\color{teal}\circled{$4$}}\ \scalebox{2}{\textcolor{red}{\fbox{\bf $12$}}}\ \scalebox{1.5}{\color{teal}\circled{$5$}}\ \scalebox{2}{\textcolor{red}{\fbox{\bf $12$}}}\ \scalebox{1.5}{\color{teal}\circled{$6$}}\ \scalebox{2}{\textcolor{red}{\fbox{\bf $12$}}}\ \scalebox{1.5}{\color{teal}\circled{$7$}}
			\end{center}
			Khi đó giữa $6$ bạn khối $12$ tạo thành $7$ khoảng trống, để đảm bảo không có $2$ học sinh cùng  khối đứng cạnh nhau, có $2$ trường hợp sau:
			\begin{enumerate}
				\item [TH1.] Xếp $6$ bạn khối $10$, $11$ vào các vị trí {\color{teal}\circled{$1$}}\ {\color{teal}\circled{$2$}}\ {\color{teal}\circled{$3$}}\  {\color{teal}\circled{$4$}}\ {\color{teal}\circled{$5$}}\ {\color{teal}\circled{$6$}} hoặc {\color{teal}\circled{$2$}}\ {\color{teal}\circled{$3$}}\ {\color{teal}\circled{$4$}}\ {\color{teal}\circled{$5$}}\ {\color{teal}\circled{$6$}}\ {\color{teal}\circled{$7$}}, có $2\cdot 6!$ cách.
				\item [TH2.] Chọn $1$ cặp $10-11$ đứng cạnh nhau, cùng với $4$ bạn  còn lại xếp vào các vị trí {\color{teal}\circled{$2$}}\ {\color{teal}\circled{$3$}}\ {\color{teal}\circled{$4$}}\ {\color{teal}\circled{$5$}}\ {\color{teal}\circled{$6$}}, có $\mathrm{C}_4^1\cdot \mathrm{C}_2^1\cdot 2!\cdot 5!$.
			\end{enumerate}
			 Vậy có $6!\left(2\cdot 6!+\mathrm{C}_4^1\cdot \mathrm{C}_2^1\cdot 2!\cdot 5!\right)=2\ 419\ 200$ (cách).
		\end{itemchoice}
	 }
\end{ex}


%%%%%%%%%%%%%%%% TRẢ LỜI NGẮN %%%%%%%%%%%%%%%%

% nguồn: MaHoaNH_TraLoiNgan_Lop10
\begin{bt}%[0D0H2-4]%[ma:6c983cc0]
	Trong một cuộc tổng điều tra dân số, điều tra viên chọn ngẫu nhiên một gia đình có ba người con và quan tâm đến giới tính của ba người con này. Giả thiết rằng khả năng sinh con trai và khả năng sinh con gái là như nhau. Tính xác suất để gia đình đó có hai  con gái biết rằng gia đình đó có con gái đầu lòng.
	\shortans{0,5} 
	 \loigiai{Kí hiệu con gái là G, con trai là T.\\Kí hiệu GTT chỉ giới tính con thứ nhất là con gái, con thứ hai là con trai, con thứ ba là con trai.\\
		Tương tự như vậy, ta có không gian mẫu là tập $\Omega=\{G G G$; $G G T$; $G T  G$; $G T T\}$.\\
		Gọi $A$  là biến cố cần tìm, suy ra $ A=\{GGT$; $GTG\}$.\\
		Vậy $\mathrm{P}(A)=\dfrac{n(A)}{n(\Omega)}=\dfrac{2}{4}=0{,}5$. 
	} 
\end{bt}

% nguồn: MaHoaNH_TraLoiNgan_Lop10
\begin{bt}%[0D0H2-4]%[ma:de264974]
	Một lớp học có $15$ học sinh nam và $20$ học sinh nữ. Chọn ngẫu nhiên $10$ học sinh  để đi lao động vệ sinh. Gọi $P$ là xác suất để $10$ học sinh chọn ra đều là nam. Tính $100000 P$ (làm tròn tới hàng phần trăm)
	\shortans{1,64} 
	 \loigiai{
		Chọn ngẫu nhiên $10$ học sinh đi lao động, số phần tử của không gian mẫu là  $n(\Omega)=\mathrm{C}_{35}^{10}$.\\
		Xét biến cố $A$: \lq\lq Chọn được $10$ học  sinh đều là nam\rq\rq, số phần tử của không gian biến cố là $n(A)=\mathrm{C}_{15}^{10}$.\\
		Xác suất cần tìm là  $\mathrm{P}(A)=\dfrac{n(A)}{n(\Omega)}=\dfrac{\mathrm{C}_{15}^{10}}{\mathrm{C}_{35}^{10}}=\dfrac{1}{61132}$.\\
		Suy  ra $100000 P =\dfrac{100000}{61132}\approx 1{,}64$.
	 }
\end{bt}

% nguồn: MaHoaNH_TraLoiNgan_Lop10
\begin{bt}%[0D0H2-4]%[ma:817fa78a]
	Một tổ có $9$ học sinh nam và $3$ học sinh nữ. Chia tổ thành $3$ nhóm, mỗi nhóm $4$ người để làm $3$ nhiệm vụ khác nhau. Xác suất khi chia ngẫu  nhiên nhóm nào cũng có nữ bằng $\dfrac{a}{b}$ với $\dfrac{a}{b}$ là phân số tối giản và $a$, $b \in \mathbb{Z}$. Tính $T=a+25 b$.
	 \shortans{1,39}
	\loigiai{Số  phần tử của không gian mẫu là $n(\Omega)=\mathrm{C}_{12}^4\cdot \mathrm{C}_{8}^4 \cdot 1=34\,650$.\\
		Gọi $A$ là biến cố \lq\lq Chia mỗi nhóm có đúng một nữ và ba nam\rq\rq.\\ 
		Số cách phân chia cho nhóm $1$ là  $\mathrm{C}_3^1 \cdot \mathrm{C}_9^3=252$ (cách).\\
		Khi đó còn lại $2$ nữ $6$ nam nên số  cách phân chia cho nhóm $2$ có $\mathrm{C}_2^1 \cdot \mathrm{C}_6^3=40$ (cách).\\
		Cuối cùng còn lại bốn người thuộc về nhóm $3$ nên  có $1$ cách chọn.
		Theo quy tắc nhân ta có số phần tử của  biến cố $A$ là $n(A)=252 \cdot 40 \cdot 1=10\,080$ (cách).\\
		Vậy xác suất cần  tìm là $$\mathrm{P}(A)=\dfrac{10\,080}{34\,650}=\dfrac{16}{55}\Rightarrow\heva{&a=16 \\& b=55} \Rightarrow T=a+25 b=1\,391.$$
	} 
\end{bt}

% nguồn: MaHoaNH_TraLoiNgan_Lop10
\begin{bt}%[0D0H2-4]%[ma:c3858e24]
	Trong một lớp có $15$ học sinh nam và $10$ học sinh nữ. Giáo viên gọi ngẫu nhiên $4$ học sinh lên bảng làm bài tập. Xác suất để $4$ học sinh  được chọn có cả nam và nữ bằng bao nhiêu? (Kết quả làm tròn đến hai chữ số sau dấu phẩy).
	 \shortans{0,88}
	\loigiai{ 
		 Số phần tử của không gian mẫu $\mathrm{n}\left( \Omega\right) =\mathrm{C}_{25}^4$.\\
		Gọi $A$ là biến cố \lq\lq  $4$ học sinh được gọi đó cả nam lẫn nữ\rq\rq.\\
		Khi đó 
		$$\mathrm{n} (A)=\mathrm{C}_{15}^1 \mathrm{C}_{10}^3+\mathrm{C}_{15}^2  \mathrm{C}_{10}^2+\mathrm{C}_{15}^3 \mathrm{C}_{10}^1.$$
		Vậy  xác suất cần tìm là
		 $$\mathrm{P}(A)=\dfrac{\mathrm{n} (A)}{\mathrm{n} (\Omega)}=\dfrac{\mathrm{C}_{15}^1 \mathrm{C}_{10}^3+\mathrm{C}_{15}^2 \mathrm{C}_{10}^2+\mathrm{C}_{15}^3 \mathrm{C}_{10}^1}{\mathrm{C}_{25}^4}=\dfrac{443}{506} \approx 0{,}88.$$
	} 
\end{bt}

% nguồn: MaHoaNH_TraLoiNgan_Lop10
\begin{bt}%[0D0H2-4]%[ma:c980ee88]
	Tại buổi biểu diễn văn nghệ \lq\lq CHILL COVER \rq\rq\, của trường THPT Lê Lợi tổ chức có sự tham gia của $5$ tiết mục văn nghệ của khối 12, $4$ tiết mục văn  nghệ của khối 11 và $4$ tiết mục văn nghệ của khối 10. Tính xác suất để ban tổ chức bố trí các tiết mục biểu diễn sao cho không có hai tiết mục văn nghệ nào của khối 12 biểu diễn liên tiếp nhau. (làm tròn đến hàng phần chục)
	\shortans{0,1} 
	 \loigiai{
		Để bố trí các tiết mục biểu diễn, số phần tử của  không gian mẫu là $n(\Omega)=13!$.\\
		Xét  biến cố không có hai tiết mục nào của khối 12 biểu diễn liên tiếp nhau.
		\begin{itemize}
			\item Xếp $8$ tiết mục của hai khối  10 và 11 liên tiếp nhau, ta có $8!$ cách xếp.
			\item Để không có tiết mục nào của khối 12 liên  tiếp nhau thì các tiết mục đó được xếp vào các khoảng hở của các tiết mục ở hai khối còn lại, hoặc biểu diễn đầu hoặc cuối, có $\mathrm{A}_9^5$ cách.
		\end{itemize}
		Vậy số khả năng thuận lợi cho biến  cố $A$ là $n(A)=8!\cdot \mathrm{A}_9^5$.\\
		Xác suất cần tìm  là $\mathrm{P}(A)=\dfrac{n(A)}{n(\Omega)}=\dfrac{8!\cdot \mathrm{A}_9^5}{13!}=\dfrac{14}{143}\approx 0{,}1$.
	 }
\end{bt}

% nguồn: MaHoaNH_TraLoiNgan_Lop10
\begin{bt}%[0D0H2-4]%[ma:1dcd5b34]
	Có $13$ học sinh của một trường THPT đạt danh hiệu học sinh xuất sắc trong đó khối $12$ có $8$ học sinh nam và $3$ học sinh nữ, khối $11$ có $2$ học sinh  nam. Chọn ngẫu nhiên $3$ học sinh bất kỳ để trao thưởng,
	xác suất để $3$ học sinh được chọn có cả nam và  nữ đồng thời có cả khối $11$ và khối $12$ bằng $\dfrac{a}{b}$ với $\dfrac{a}{b}$ là phân số tối giản và $a$, $b \in \mathbb{Z}$. Tính $ T=3 a+2 b$.
	\shortans{743} 
	\loigiai{Không gian mẫu là số  cách chọn ngẫu nhiên $3$ học sinh từ $13$ học sinh.\\
		Suy ra số phần tử của không gian mẫu là $n(\Omega)=\mathrm{C}_{13}^3=286$. \\Gọi  $A$ là biến cố \lq\lq $3$ học sinh được chọn có cả nam và nữ đồng thời có cả khối $11$ và khối $12$\rq\rq .\\ Ta xét các trường hợp sau:
		\begin{itemize}
			\item \textbf{Trường hợp 1.} Chọn $1$ học sinh  khối $11$, $1$ học sinh nam khối $12$ và $1$ học sinh nữ khối $12$ nên có $\mathrm{C}_2^1\cdot \mathrm{C}_8^1\cdot \mathrm{C}_3^1=48$ cách.
			\item \textbf{Trường hợp 2.} Chọn $1$  học sinh khối $11$, $2$ học sinh nữ khối $12$ có $\mathrm{C}_2^1 \cdot \mathrm{C}_3^2=6$ cách.
			\item \textbf{Trường hợp 3.} Chọn $2$ học sinh khối $11$,  $1$ học sinh nữ khối $12$ có $\mathrm{C}_2^2 \cdot \mathrm{C}_3^1=3$ cách.
		\end{itemize}
		Suy ra số phần  tử của biến cố $A$ là $n(A)=48+6+3=57$.\\
		Vậy  xác suất cần tính $$\mathrm{P}(A)=\dfrac{n(A)}{n(\Omega)}=\dfrac{57}{286}\Rightarrow \heva{&a=57 \\& b=286} \Rightarrow T=3 a+2 b=743.$$
	} 
\end{bt}

% nguồn: MaHoaNH_TraLoiNgan_Lop10
\begin{bt}%[0D0V2-4]%[ma:3ed25ca9]
	Lớp $10A8$ có $10$ bạn nữ và $20$ bạn nam. Thầy giáo chủ nhiệm cần chọn ra $5$ bạn để biên tập Video giới thiệu về lớp mình. Tính xác suất để  thầy có thể chọn được $2$ bạn nam và $3$ bạn nữ (kết quả làm tròn đến hàng phần trăm). \shortans{0,16}
	\loigiai{ 
		Ta có $n(\Omega) = \mathrm{C}_{30}^{5}$.\\ 
		Gọi $A\colon$ \lq\lq Chọn được $2$ bạn nam và $3$ bạn  nữ\rq\rq.\\
		 Ta có $n(A) = \mathrm{C}_{20}^{2}\cdot \mathrm{C}_{10}^{3}$.\\
		Xác suất để thầy chủ nhiệm có thể chọn  được 2 bạn nam và 3 bạn nữ là
		$$ \mathrm{P}(A) = \dfrac{n(A)}{n(\Omega)} = \dfrac{\mathrm{C}_{20}^{2}\cdot \mathrm{C}_{10}^{3}}{\mathrm{C}_{30}^{5}}  = \dfrac{3800}{23751} \approx 0{,}16. $$
	} 
\end{bt}

% nguồn: MaHoaNH_TraLoiNgan_Lop10
\begin{bt}%[0D0V2-4]%[ma:6ea840b8]
	Lớp  A có $5$ học sinh xuất sắc, lớp B có $10$ học sinh xuất sắc. Chọn ngẫu nhiên $2$ trong các học sinh đó đi dự lễ tuyên dương học sinh tiêu biểu. Tính xác suất để $2$ học sinh được chọn từ cùng một lớp (kết quả làm tròn đến hàng phần trăm). \shortans{0,52}
	\loigiai{ 
		Số phần tử của không gian mẫu là $n(\Omega) =  \mathrm{C}_{15}^{2}$.\\
		 Gọi $A\colon$ \lq\lq $2$ học sinh được chọn từ cùng một lớp \rq\rq.\\
		Khi đó $n(A) = \mathrm{C}_{5}^{2} +  \mathrm{C}_{10}^{2}$.\\
		Xác suất để $2$ học sinh được  chọn từ cùng một lớp là
		$$ \mathrm{P}(A)  = \dfrac{n(A)}{n(\Omega)} = \dfrac{\mathrm{C}_{5}^{2} + \mathrm{C}_{10}^{2}}{\mathrm{C}_{15}^{2}} = \dfrac{11}{21} \approx 0{,}52. $$
	} 
\end{bt}

% nguồn: MaHoaNH_TraLoiNgan_Lop10
\begin{bt}%[0D0V2-4]%[ma:b05393a4]
	Lớp $10A8$ có $10$ bạn nữ và $20$ bạn nam. Thầy giáo chủ nhiệm cần chọn ra $5$ bạn để biên tập Video giới thiệu về lớp mình. Tính xác suất để thầy có thể chọn nhiều nhất $1$  bạn nam (kết quả làm tròn đến hàng phần trăm). \shortans{0,03}
	 \loigiai{
		Ta có $n(\Omega) = \mathrm{C}_{30}^{5}$.\\ 
		Gọi $A\colon$ \lq\lq Chọn được nhiều nhất $1$ bạn  nam\rq\rq.\\
		Ta có $n(A) = \mathrm{C}_{20}^{0}\cdot \mathrm{C}_{10}^{5} + \mathrm{C}_{20}^{1}\cdot  \mathrm{C}_{10}^{4}$.\\
		Xác  suất để thầy chủ nhiệm có thể chọn được 2 bạn nam và 3 bạn nữ là
		$$ \mathrm{P}(A) = \dfrac{n(A)}{n(\Omega)} = \dfrac{\mathrm{C}_{20}^{0}\cdot \mathrm{C}_{10}^{5} + \mathrm{C}_{20}^{1}\cdot \mathrm{C}_{10}^{4}}{\mathrm{C}_{30}^{5}} = \dfrac{106}{3393} \approx 0{,}03.  $$
	 }
\end{bt}

% nguồn: MaHoaNH_TraLoiNgan_Lop10
\begin{bt}%[0D0V2-4]%[ma:e2c8b658]
	Xếp ngẫu nhiên $4$ bạn nam và $4$ bạn nữ thành một hàng dọc. Tính xác suất của biến cố \lq\lq Xếp được các bạn nam và bạn nữ đứng xen kẽ nhau\rq\rq\,  (kết quả làm tròn đến hàng phần trăm).
	\shortans  {$0{,}03$}
	 \loigiai{
		Giả sử các vị trí của hàng dọc được đánh số thứ tự từ  đầu hàng là $1,2,3,4,5,6,7,8$.\\
		Số cách xếp 8  bạn thành một hàng dọc là $8!=40320$.\\
		Xếp các bạn nam và bạn nữ đứng xen kẽ nhau có  hai trường hợp
		\begin{itemize}
			\item {\bf Trường hợp 1.} Các bạn nam đứng ở các vị trí số lẻ còn các bạn nữ đứng ở các vị trí số chẵn.\\ Số cách xếp  như vậy là $4!\cdot4!=576$.
			\item {\bf Trường hợp 2.} Các bạn nữ đứng ở các vị trí số lẻ còn các bạn nam đứng ở các vị trí số chẵn.\\ Số cách xếp như vậy  là $4!\cdot 4!=576$.
		\end{itemize}
	Vậy xác suất của biến cố \lq\lq Xếp được các bạn nam và bạn nữ đứng xen kẽ nhau\rq\rq\,  là $\dfrac{576+576}{40320}=\dfrac{1}{35}\approx 0{,}03$.}
\end{bt}

% nguồn: MaHoaNH_TraLoiNgan_Lop10
\begin{bt}%[0D0V2-4]%[ma:78f54449]
	Có $5$ học sinh nam và $10$ học sinh nữ, trong các học sinh nữ có Vy và Quyên, Lan. Xếp những học sinh này thành một hàng ngang. Xác suất để mỗi bạn nam đều đứng giữa hai bạn nữ đồng thời Vy, Quyên,  Lan đứng cạnh nhau bằng $\dfrac{a}{b}$ với $\dfrac{a}{b}$ là phân số tối giản và $a$, $b \in \mathbb{Z}$. Tính giá trị biểu thức $T=b-2\,025a$.
	 \shortans{0,14}
	\loigiai{ 
		Xét phép thử: \lq\lq Xếp $5$ học sinh nam và $10$  học sinh nữ thành một hàng ngang\rq\rq, suy ra $n(\Omega)=15!$.\\
		Gọi biến cố $X$:  \lq\lq Mỗi bạn nam đều đứng giữa hai bạn nữ đồng thời Vy, Quyên, Lan luôn đứng cạnh nhau\rq\rq .
		\begin{itemize}
			\item  \textbf{Bước 1.} Xếp Vy, Quyên, Lan đứng cạnh nhau có $3!$ cách.
			\item \textbf{Bước 2.} Xếp Vy, Quyên, Lan  và $7$ bạn còn lại vào $8$ vị trí có $8 !$ cách.
			\item \textbf{Bước 3.} Chọn $5$ khoảng trống trong $7$  khoảng trống giữa $8$ vị trí ở bước $2$ cho $5$ bạn nam có $A_7^5$ cách
		\end{itemize}
		 Suy ra $n(X)=3!\cdot 8!\cdot \mathrm{A}_7^5$.\\
		Vậy xác suất của biến cố $X$  là $$\mathrm{P}(X)=\dfrac{3!8!\mathrm{A}_7^5}{15!}=\dfrac{1}{2\,145}\Rightarrow\heva{&a=1 \\& b=2\,145} \Rightarrow T=b-2\,025 a=120.$$
	} 
\end{bt}

% nguồn: MaHoaNH_TraLoiNgan_Lop10
\begin{bt}%[0D0V2-4]%[ma:75e83513]
	Lớp 11A6 có $45$ học sinh gồm $20$ nam và $25$ nữ. Chọn  ngẫu nhiên $8$ bạn đi học đối tượng Đoàn, Tính xác suất để $8$ bạn được chọn có cả nam và nữ (làm tròn kết quả đến chữ số thập phân thứ hai).
	\shortans{0,99} 
	 \loigiai{
		Ta có $n(\Omega) = \mathrm{C}_{45}^8$.\\ 
		Gọi $A$ là biến cố \lq\lq chọn ra $8$ bạn đi học đối tượng Đoàn sao cho trong đó có cả nam lẫn nữ \rq\rq.  \\
		Suy  ra $\overline{A}$ là biến cố \lq\lq chọn ra $8$ bạn đi học đối tượng Đoàn sao cho trong đó chỉ có nam hoặc nữ\rq\rq.\\
		Chọn ra $8$ bạn học sinh bất kì trong số  $45$ học sinh có $\mathrm{C}_{45}^8$ cách.\\
		Chọn ra $8$ bạn học sinh bất kì trong số  $20$ học sinh nam có $\mathrm{C}_{20}^8$ cách.\\
		Chọn ra  $8$ bạn học sinh bất kì trong số $25$ học sinh nữ có $\mathrm{C}_{25}^8$ cách.\\
		Vậy số cách chọn  ra $8$ bạn đi học đối tượng Đoàn sao cho trong đó có cả nam lẫn nữ là
		$$\mathrm{C}_{45}^8-\mathrm{C}_{20}^8-\mathrm{C}_{25}^8=214\,345\,650\;  \text{cách}.$$
		Xác suất  là $P(A) =\dfrac{n(A)}{n(\Omega)}= \dfrac{214\,345\,650}{\mathrm{C}_{45}^8} \approx 0{,}99$.
	} 
\end{bt}


%%%%%%%%%%%%%%%% TỰ LUẬN %%%%%%%%%%%%%%%%

% nguồn: MaHoaNH_TuLuan_Lop10
\begin{bt}%[0D0H2-4]%[ma:2057c2c4]
	Một nhóm học sinh gồm có $9$ nam, $3$  nữ. Tính xác suất để khi chọn ngẫu nhiên $4$ người thì có đúng $1$ nữ.
	\loigiai{Số phần tử của không gian  mẫu là $n(\Omega) = \mathrm{C}_{12}^{4}=495$.\\
		Gọi $A$ là biến cố: \lq \lq Trong $4$ người được chọn thì có đúng  $1$ nữ \rq \rq. \\
		Để chọn $4$ người có đúng $1$ nữ thì phải  chọn $1$ nữ và $3$ nam có $\mathrm{C}_{3}^{1} \cdot \mathrm{C}_{9}^{3}=252$ cách.\\ Do đó $n(A)=252$.\\
		Vậy xác suất  cần tìm là $\mathrm{P}(A)=\dfrac{n(A)}{n(\Omega)}=\dfrac{252}{495}=\dfrac{28}{55}$.
	 }
\end{bt}

% nguồn: MaHoaNH_TuLuan_Lop10
\begin{bt}%[0D0H2-4]%[ma:6e38227c]
	Một đội gồm $5$ nam và $8$ nữ. Lập một nhóm gồm $4$ người hát tốp ca, tính xác suất để trong $4$ người được chọn có ít nhất $3$  nữ.
	\loigiai{Chọn tùy ý $4$ người từ $13$ người nên số  phần tử của không gian mẫu là $n(\Omega)=\mathrm{C}_{30}^{3}$.\\
		Gọi $A$  là biến cố \lq \lq Trong $4$ người được chọn có ít nhất $3$ nữ\rq \rq. \\
		Ta có hai  trường hợp như sau:\\
		\textbf{Trường hợp 1:} Chọn $3$ nữ và $1$ nam, có $\mathrm{C}_{8}^{3} \cdot \mathrm{C}_{5}^{1}$  cách.\\
		\textbf{Trường hợp 2:}  Chọn cả $4$ nữ, có $\mathrm{C}_{8}^{4}$ cách.\\
		Suy ra số phần tử của biến cố $A$ là $n(A)=\mathrm{C}_{8}^{3} \cdot  \mathrm{C}_{5}^{1}+\mathrm{C}_{8}^{4}=350$.\\
	Vậy xác suất cần tính là  $\mathrm{P}(A)=\dfrac{n(A)}{n(\Omega)}=\dfrac{350}{715}=\dfrac{70}{143}$. }
\end{bt}
